\documentclass{ieeetj}
\usepackage{cite}
\usepackage{amsmath,amssymb,amsfonts}
\usepackage{algorithmic}
\usepackage{graphicx,color}
\usepackage{textcomp}
\usepackage{xcolor}
\usepackage{hyperref}
\usepackage{multirow}
\usepackage{booktabs}
\usepackage{algorithm}
\usepackage{algpseudocode}
\usepackage{adjustbox}
\hypersetup{hidelinks=true}
\graphicspath{{assets/}{report/images/}}

\def\BibTeX{{\rm B\kern-.05em{\sc i\kern-.025em b}\kern-.08em
    T\kern-.1667em\lower.7ex\hbox{E}\kern-.125emX}}
\AtBeginDocument{\definecolor{tmlcncolor}{cmyk}{0.93,0.59,0.15,0.02}\definecolor{NavyBlue}{RGB}{0,86,125}}

\def\authorrefmark#1{\ensuremath{^{\textbf{#1}}}}

\begin{document}
\receiveddate{XX Month, XXXX}
\reviseddate{XX Month, XXXX}
\accepteddate{XX Month, XXXX}
\publisheddate{XX Month, XXXX}
\currentdate{XX Month, XXXX}
\doiinfo{XXXX.2024.XXXXXX}

\markboth{Girish et al.}{Explainable Deep Learning for Sleep Disorder Detection}

\title{Explainable Deep Learning Approach for Sleep Disorder Detection and Personalized Lifestyle Recommendations}

\author{Girish Sundararajan\authorrefmark{1}, Harishankar Binu Nair\authorrefmark{1}, and Simi Surendran\authorrefmark{2}}
\affil{Department of Artificial Intelligence, Amrita School of Computing, Amrita Vishwa Vidyapeetham, Amritapuri, Kerala 690525, India}
\affil{Department of Artificial Intelligence, Amrita School of Computing, Amrita Vishwa Vidyapeetham, Amritapuri, Kerala 690525, India}
\corresp{Corresponding author: Girish Sundararajan (email: giri090205@gmail.com).}

\begin{abstract}
Sleep disorders affect a substantial portion of the global population, yet current diagnostic systems fragment care across separate staging and disorder detection pipelines with limited interpretability. Existing deep learning approaches achieve high accuracy but function as black boxes, preventing clinical adoption and patient engagement. This study addresses the need for a unified explainable system supporting both clinical decision making and patient self management. An explainable deep learning framework was developed providing unified sleep architecture analysis and disorder risk assessment. The system combines a temporal spectral fusion model for sleep staging with multi-label classifiers for apnea, insomnia, and restless legs syndrome detection. Explainability methods operate on two levels. For clinicians, attention rollout provides per epoch temporal influence scores, GradCAM localizes signal regions driving predictions, and Monte Carlo dropout quantifies prediction uncertainty to flag cases requiring review. For patients, risk scores present disorder probabilities as percentages, feature importance highlights key contributing factors, counterfactual explanations suggest actionable changes, and a conversational interface answers questions in plain language. The unified platform was evaluated on the Multi-Ethnic Study of Atherosclerosis dataset containing over two thousand polysomnography recordings. Sleep staging achieved macro F1 score of 0.788 with accurate detection of wake and deep sleep stages. Disorder detection achieved F1 scores of 0.784 for apnea, 0.388 for insomnia, and 0.182 for restless legs syndrome. ElasticNet stability selection improved insomnia detection by 35.9 percent over standard regularization by preserving complementary feature groups. The explainability pipeline successfully identified uncertain predictions requiring clinician attention and generated personalized recommendations actionable by patients. Analysis confirmed that self reported disorder labels align with behavioral features while periodic limb movement markers provide direct physiological signatures for restless legs syndrome. The proposed framework demonstrates that a unified explainable approach can bridge the gap between accurate automated analysis and clinical utility.
\end{abstract}

\begin{IEEEkeywords}
Sleep disorder, sleep staging, explainable AI, deep learning, polysomnography, transformer, temporal-spectral fusion
\end{IEEEkeywords}

\maketitle

\section{INTRODUCTION}
\IEEEPARstart{S}{leep} health is foundational to cognitive performance, metabolic regulation, and cardiovascular stability \cite{berry2012aasm, berry2020aasm}. Disrupted sleep is linked to a wide range of chronic conditions, underscoring the need for accurate diagnostic tools.

Sleep architecture is composed of five stages—wake, N1, N2, N3, and REM—each distinguished by characteristic electroencephalogram (EEG) patterns \cite{supratak2017deepsleepnet, phan2022sleeptransformer}. Reliable identification of these stages is essential for assessing overall sleep quality and for downstream clinical analyses.

Lifestyle factors such as caffeine consumption, shift work, and physical activity significantly influence sleep architecture and modulate the risk of sleep‑related disorders \cite{fu2023tsanet, kim2018apna}. These influences are reflected in both temporal dynamics and spectral characteristics of the EEG signal.

The three most prevalent sleep disorders addressed in this study are obstructive sleep apnea (OSA), insomnia, and restless‑legs syndrome (RLS). OSA is characterised by recurrent upper‑airway collapse during sleep, insomnia reflects hyper‑arousal and difficulty maintaining sleep, and RLS is associated with dopaminergic dysfunction leading to uncomfortable leg sensations.

Epidemiological data indicate that OSA affects approximately 10–17\% of adults, insomnia 10–30\%, and RLS 5–10\%, imposing substantial health and economic burdens \cite{rashid2022osa, karthikeyan2023insomnet, seifollahi2023osarls}. These prevalence rates highlight the public health importance of effective detection and management strategies.

Polysomnography (PSG) provides the gold‑standard measurement for sleep staging and disorder diagnosis by recording EEG, EOG, EMG, and respiratory signals. However, manual scoring of PSG recordings is time‑intensive and subject to inter‑rater variability \cite{berry2020aasm, jha2024slumbernet}.

Recent automation trends have shifted toward deep‑learning approaches, particularly temporal‑spectral fusion transformers, which achieve near‑human staging accuracy while enabling interpretability through attention rollout, Grad‑CAM visualisations, and Monte‑Carlo dropout uncertainty estimation \cite{cong2025bits, sun2025adg, phan2022sleeptransformer}.

In this work we present a unified, explainable framework that (i) fuses raw EEG and spectral features to achieve robust sleep staging, (ii) detects OSA, insomnia, and RLS using a combination of questionnaire and physiological features enhanced by ElasticNet‑based bootstrap stability selection, and (iii) delivers clinician‑focused attention‑based explanations alongside patient‑focused risk scores, feature importance visualisations, and actionable counterfactual guidance. Figure~\ref{fig:overview} illustrates the overall system architecture.

\begin{figure*}
    \centering
    \includegraphics[width=1\linewidth]{assets/prj 48 methodology.jpg}
    \caption{Overview of the 3‑tier system architecture proposed in this work.}
    \label{fig:overview}
\end{figure*}

The contributions of this work are as follows. First, a temporal‑spectral fusion architecture combines multi‑scale dilated convolutions with spectral power band features, achieving state‑of‑the‑art performance on N1 detection via concatenation fusion. Second, ElasticNet bootstrap stability selection captures complementary feature groups, improving insomnia detection by 35.9\% over L1 regularisation. Third, an explainability pipeline offers clinician‑facing attention rollout, GradCAM on the AdaptiveAtrousPyramid projection layer, and MC‑Dropout uncertainty with dataset‑calibrated confidence tiers, alongside patient‑facing SHAP feature importance, counterfactual lifestyle targets, and a Groq Cloud LLM advisor. Fourth, we demonstrate that self‑reported disorder labels align with behavioural questionnaire features rather than raw physiological signals, informing future diagnostic system design.
\section{RELATED WORKS}

Early automated sleep staging systems relied on handcrafted feature extraction paired with traditional classifiers. Guo et al.~\cite{guo2017hilbert} proposed a method using Hilbert-Huang transform and sample entropy features with SVM classifiers, while Hassan et al.~\cite{hassan2017decision} utilized tunable Q-factor wavelet transform with bootstrap aggregating for multi-stage classification. These approaches typically achieved accuracies ranging from 75\% to 85\% on standard benchmarks but struggled with ambiguous stages such as N1. The fundamental limitation of such methods lay in their dependence on human-engineered features, which required substantial domain expertise and often failed to capture complex physiological patterns.

The introduction of deep learning transformed the field by enabling end-to-end feature learning from raw signals. Supratak et al.~\cite{supratak2017deepsleepnet} proposed DeepSleepNet, which combined CNN and BiLSTM architectures to extract time-invariant features while learning transition rules between sleep stages. Phan et al.~\cite{phan2018deepsleepnet} extended this with SeqSleepNet, a hierarchical recurrent neural network operating on sequences of epochs rather than individual samples. While these methods improved accuracy, they functioned as black boxes, providing no mechanism for clinicians to understand model predictions.

Residual architectures have further advanced performance. Jha et al.~\cite{jha2024slumbernet} proposed SlumberNet based on ResNet architecture, achieving 97\% accuracy on mouse sleep staging through skip connections that preserved gradient flow. However, the opacity of such deep networks created skepticism among clinicians, which Phan et al.~\cite{phan2022sleeptransformer} identified as a primary barrier to clinical adoption. They argued that two capabilities were essential for practical deployment: interpretability explaining how decisions were made, and uncertainty quantification enabling deferral of low-confidence predictions to human experts.

Contemporary approaches have increasingly recognized that sleep stages exhibit distinguishing characteristics in both temporal and spectral domains. Fu et al.~\cite{fu2023tsanet} proposed TSA-Net, a two-stream architecture extracting features simultaneously from raw EEG waveforms and spectral representations, followed by multi-head self-attention and conditional random fields for transition rule modeling. Their approach achieved 86.6\% accuracy on Sleep-EDF-20 and 82.2\% on Sleep-EDF-78, demonstrating that temporal-spectral fusion consistently outperformed single-modality approaches. Cong et al.~\cite{cong2025bits} advanced this paradigm with BiTS-SleepNet, introducing cross-attention fusion that dynamically weighted temporal and spectral contributions per epoch. Their two-stage architecture combining dual-stream feature extraction with Bi-GRU context learning and CRF transition rules reached 88.5\% accuracy on Sleep-EDF-20.

Sun and Zhao~\cite{sun2025adg} proposed ADG-SleepNet, combining multi-scale dilated temporal convolutions with adaptive gating and squeeze-excitation attention, achieving 87.1\% accuracy while maintaining computational efficiency. Huang et al.~\cite{huang2025mhas} introduced MHASSNet, which integrated multi-scale convolution, hybrid attention mechanisms, and state space models for enhanced temporal contextual understanding. These works established hybrid CNN-attention architectures as the dominant paradigm for high-performance sleep staging.

While accuracy improved, the interpretability challenge remained largely unaddressed until recent years. Phan et al.~\cite{phan2022sleeptransformer} proposed SleepTransformer, the first pure transformer-based sleep staging model, which leveraged self-attention scores for interpretability at both epoch and sequence levels. Their entropy-based uncertainty quantification enabled deferral of ambiguous epochs to human experts, with experiments showing that approximately 90\% of epochs per night could be confidently automated. Lee et al.~\cite{lee2025sleepxvit} extended explainability to multimodal PSG inputs with SleepXViT, using layer-wise relevance propagation to generate pixel-level heatmaps aligned with AASM clinical guidelines. Their inter-epoch Vision Transformer achieved 81.9\% macro F1 on the large-scale KISS dataset while providing clinically interpretable visualizations.

While sleep staging has received substantial attention, automated detection of sleep disorders from PSG recordings remains comparatively underexplored. Rashid et al.~\cite{rashid2022osa} proposed a lightweight CNN for OSA detection using single-lead ECG, achieving 94.7\% per-recording accuracy through heart rate variability and ECG-derived respiration features. Seifollahi et al.~\cite{seifollahi2023osarls} addressed co-morbid OSA and RLS detection using Synchrosqueezing Wavelet Transform of ECG and EMG signals, with multimodal fusion achieving 94.1\% accuracy compared to 89.2\% for ECG alone. Their work demonstrated that multi-modal fusion consistently outperformed single-modality approaches for disorder detection.

Insomnia detection presents distinct challenges due to its cognitive and behavioral components. Karthikeyan et al.~\cite{karthikeyan2023insomnet} proposed InsomNet, utilizing continuous wavelet transform of ECG signals with ResNet-LSTM architectures to achieve 93.8\% accuracy on hospital datasets. However, they noted that questionnaire-derived insomnia labels created fundamental alignment gaps with physiological signals, limiting the achievable performance from cardiac features alone.

Few studies have attempted multi-disorder classification within unified frameworks. Islam et al.~\cite{islam2024multishel} proposed a Stacked Heterogeneous Ensemble Learning approach combining CNN, BiLSTM, and XGBoost classifiers with SHAP explainability, achieving 96.3\% accuracy for apnea and 94.8\% for insomnia detection. Kumar et al.~\cite{kumar2021lstmcnn} demonstrated LSTM-CNN hybrid architectures for four-disorder classification, attaining 88.5\% overall accuracy using multi-modal PSG inputs. However, these ensemble and hybrid approaches incurred substantial computational costs and lacked integrated uncertainty quantification suitable for clinical workflows.

Several persistent limitations constrain the clinical utility of existing approaches. Most deep learning models are trained and evaluated on clean laboratory PSG recordings, with limited cross-dataset validation. Disorder detection studies predominantly focus on single conditions, particularly OSA, while clinical reality involves significant co-morbidity between apnea, insomnia, and restless legs syndrome. Explainability methods remain limited to attention visualization or post-hoc analysis without actionable uncertainty quantification. Perhaps most critically, few studies integrate sleep staging with disorder detection within unified frameworks, though sleep architecture abnormalities serve as primary diagnostic indicators for multiple disorders.

This work addresses these limitations through a unified temporal-spectral fusion architecture integrating multi-scale dilated convolutions with concatenation-based feature fusion, which outperformed bidirectional cross-attention in ablation experiments. Unlike existing disorder-specific models, the proposed framework simultaneously performs sleep staging and multi-disorder detection, capturing interdependencies between sleep architecture and pathology. The comprehensive explainability pipeline combines GradCAM for epoch-level signal localization, attention rollout for sequence-level context visualization, and Monte Carlo dropout for predictive uncertainty quantification, providing clinicians with actionable insights at multiple temporal scales.

% \section{METHODS}

% The proposed framework integrates four complementary modules into a unified end-to-end pipeline for automated sleep analysis, disorder detection, and personalised explanation, as illustrated in Figure~\ref{fig:methodology}. The first module covers dataset acquisition and description, detailing the two large-scale publicly available polysomnography datasets employed in this work. The second module describes data preprocessing and signal representation, encompassing channel selection, signal conditioning, epoch segmentation, dual-stream temporal and spectral feature engineering, and the construction of subject-level clinical feature vectors used for downstream disorder detection. The third module presents the sleep staging architecture, a temporal-spectral fusion transformer that learns discriminative representations from raw EEG waveforms and engineered spectral features simultaneously, models long-range sequential dependencies across epochs through a transformer encoder, and classifies each epoch into one of five sleep stages. The fourth module details the sleep disorder detection pipeline, a multi-label classification system that aggregates full-night clinical features to identify obstructive sleep apnea, insomnia, and restless legs syndrome through an ensemble of interpretable machine learning models. The fifth module describes the explainability system, which produces two distinct tiers of interpretive output: clinical-grade tools enabling sleep technologists and physicians to audit model predictions at the signal level, and patient-facing explanations that translate probabilistic outputs into accessible risk summaries, actionable guidance, and conversational support.


% \subsection{Dataset Description}

% Automated sleep staging has benefited from a growing collection of publicly available polysomnography corpora. Seminal benchmarks include the Sleep Cassette and Sleep Telemetry recordings released by Kemp et al.~\cite{sleepedfx}, the MASS dataset compiled from multiple Montreal laboratories, the SHHS cohort linking sleep characteristics to cardiovascular risk, and the CHAT and MrOS studies targeting paediatric and older male populations respectively. More recent additions include the CFS family study and large multicentre resources distributed through the National Sleep Research Resource~\cite{nsrr}. These datasets vary substantially in acquisition hardware, electrode placement, scoring convention, and demographic composition, producing a heterogeneous benchmark landscape that renders cross-dataset generalisation a persistent challenge. This study employs two of these publicly available resources: the Sleep-EDF Expanded dataset for supervised training and validation of the sleep staging model, and the Multi-Ethnic Study of Atherosclerosis (MESA) Sleep dataset for sleep disorder association analysis and cross-dataset generalisation evaluation. Together they enable assessment of both staging performance on a well-characterised benchmark and clinical relevance on a large population-representative cohort with adjudicated disorder labels.

% \subsubsection{Sleep-EDF Expanded Dataset}

% The Sleep-EDF Expanded dataset~\cite{sleepedfx} was obtained from PhysioNet~\cite{physiobank}, originally published in 2018 as version 1.0.0. The dataset consists of whole-night PSG recordings collected as part of two longitudinal studies: Sleep-Cassette (SC) and Sleep-Telemetry (ST). The SC subset includes recordings from healthy subjects without sleep-related medication, while the ST subset contains recordings from subjects with difficulty initiating sleep, with nights recorded under both placebo and temazepam conditions.

% Across both subsets, the dataset comprises recordings from over 100 Caucasian adult subjects, spanning an age range of approximately 25 to 101 years, with each subject contributing one or more overnight recordings. In total, the expanded dataset contains 197 recorded nights, with each night lasting approximately 8 hours.

% PSG data were acquired using standard laboratory-grade equipment, including EEG (Fpz--Cz and Pz--Oz electrode derivations), electrooculogram (EOG), and chin electromyogram (EMG) signals. Both EEG and EOG channels were sampled at 100~Hz. In this work, only the Fpz--Cz EEG channel was used to maintain consistency across datasets and to enable single-channel sleep staging.

% Sleep stages were manually annotated by trained sleep experts following the Rechtschaffen and Kales (R\&K) standard. Each recording was segmented into 30-second epochs, with labels assigned to one of the following stages: Wake (W), N1, N2, N3, N4, REM, Movement Time, and Unknown. In accordance with current AASM conventions, N4 epochs were merged into N3, and non-physiological stages (Movement Time and Unknown) were excluded during preprocessing.

% The resulting class distribution (in number of 30-second epochs) used in this study is as shown in Figure~\ref{fig:stagedist}. This distribution reflects the natural imbalance present in nocturnal sleep architecture and motivates the use of deep learning-based models capable of learning better representations under severe class imbalance.

% \begin{figure}
%     \centering
%     \includegraphics[width=1\linewidth]{assets/stages uncleaned.png}
%     \caption{Epoch-wise natural Sleep Stage Distribution from Sleep-EDFx dataset.}
%     \label{fig:stagedist}
% \end{figure}

% \subsubsection{MESA Sleep Dataset}

% The MESA Sleep dataset~\cite{mesa} was accessed through the National Sleep Research Resource (NSRR)~\cite{nsrr} and was released in 2016 as part of the Multi-Ethnic Study of Atherosclerosis. Unlike Sleep-EDF, MESA is a large-scale, population-based cohort designed to investigate the relationship between sleep characteristics and cardiometabolic health outcomes.

% The dataset includes over 2,200 middle-aged and older adults (aged approximately 45--84 years) from diverse ethnic backgrounds, with each participant undergoing one full-night in-home PSG recording. PSG acquisition followed standardised protocols and included multi-channel EEG, EOG, EMG, respiratory airflow, oxygen saturation, and electrocardiogram (ECG) signals, sampled at frequencies ranging from 100 to 256~Hz depending on the modality.

% Sleep stages in MESA were manually scored in 30-second epochs by certified sleep technologists according to AASM guidelines. In addition to sleep staging, MESA provides clinically adjudicated labels for several common sleep disorders, making it suitable for disorder-level analysis.

% In this study, MESA was used exclusively for sleep disorder association and generalisation testing, rather than for training the sleep staging model. The disorder-wise patient distribution considered in this work is summarised in Figure~\ref{fig:disorderdist}. This complementary use of Sleep-EDF for supervised sleep staging and MESA for disorder-level evaluation enables assessment of both model generalisation, mitigating learning bias, and clinical relevance across heterogeneous recording conditions and populations.

% \begin{figure}
%     \centering
%     \includegraphics[width=1\linewidth]{assets/Disorderdist.pdf}
%     \caption{Class Distribution for Sleep Disorders from MESA Sleep dataset.}
%     \label{fig:disorderdist}
% \end{figure}

% \subsection{Data Preprocessing and Representation}

% Sleep staging research employs a broad range of physiological signals, including multi-channel EEG, EOG, chin EMG, respiratory airflow, and peripheral oxygen saturation, to discriminate between sleep stages. While multi-channel models that fuse all available modalities can achieve marginally higher accuracy on balanced benchmarks~\cite{phan2022sleeptransformer, cong2024bits}, EEG alone carries sufficient discriminative information for reliable automatic staging, as evidenced by the strong performance of single-channel systems across multiple independent evaluations~\cite{fu2023temporal}. In this work, only the Fpz-Cz EEG derivation was retained for the sleep staging model, a deliberate choice that serves two purposes: it ensures consistency between the Sleep-EDF Expanded and MESA datasets, which differ in electrode complement, amplifier hardware, and recording environment, and it reduces the computational footprint of the preprocessing pipeline without sacrificing the frequency content critical for stage discrimination. A sample 30 second window of the EEG channel is presented in Fig.~\ref{fig:EEG} with respect to different sleep stages. Raw EDF recordings were loaded using the MNE-Python library, and signals were segmented into non-overlapping 30-second epochs following the Rechtschaffen and Kales standard, yielding 3,000 samples per epoch at the native 100~Hz sampling rate. The 100~Hz sampling rate was selected as a deliberate design anchor: it satisfies the Nyquist criterion for all clinically relevant EEG frequency bands up to 45~Hz (gamma), preserves the morphological detail of transient events such as K-complexes and sleep spindles whose discriminative features reside below 20~Hz, and produces epoch tensors of tractable size for GPU training. Channels sampled at higher rates in the MESA dataset were resampled to 100~Hz using polyphase anti-alias filtering prior to segmentation. Hypnogram annotations were loaded from companion EDF files for Sleep-EDF and from NSRR-format XML files for MESA, with stage mapping following current AASM conventions: Wake, N1, N2, N3 (incorporating N4), and REM. Non-physiological labels including Movement Time and Unknown epochs were excluded from training. Recordings were temporally trimmed to retain only sleep-containing segments, discarding extended wake epochs at recording boundaries while preserving a 30-minute buffer before sleep onset and after final awakening to reduce class imbalance without distorting the natural wake-to-sleep transition distribution. A final processed stage transition across the night for a few subjects is shown in Fig.~\ref{fig:stagetransition}.

% \begin{figure}
%     \centering
%     \includegraphics[width=1\linewidth]{assets/EEG Fpz-Cz.png}
%     \caption{EEG Fpz-Cz Channel waveform comparison against different Sleep Stages.}
%     \label{fig:EEG}
% \end{figure}

% \begin{figure}
%     \centering
%     \includegraphics[width=1\linewidth]{assets/stage sequence cleaned.png}
%     \caption{Cleaned Stage Transition across the full night for few sample subjects.}
%     \label{fig:stagetransition}
% \end{figure}

% Representing each epoch by its raw waveform alone would require the model to discover frequency-domain structure from scratch through convolutional operations, a computationally redundant process given that sleep stage boundaries are well characterised by canonical frequency bands known for over four decades of research. Conversely, using only pre-computed spectral features compresses the signal into a fixed-length summary that irreversibly discards morphological transients, precisely the brief events such as sleep spindles (0.5--2 seconds), K-complexes, and REM phasic bursts that carry high discriminative value but are smoothed away by windowed spectral analysis. A dual-stream representation preserves both information types simultaneously: the temporal stream supplies the raw morphological signal, while the spectral stream provides a compact, noise-robust frequency summary that guides the model toward physiologically grounded representations. This design is supported by the ablation results of Cong et al.~\cite{cong2024bits} and Fu et al.~\cite{fu2023temporal}, who independently demonstrated 12--15 percentage point accuracy gains when combining temporal and spectral inputs over either modality alone. For the spectral stream, two complementary transform methods are employed: Discrete Wavelet Transform (DWT) using the Daubechies-4 (db4) basis at decomposition level 5, and Short-Time Fourier Transform (STFT) with a 2-second Hann window and 50 percent overlap. DWT and STFT capture different aspects of the non-stationary EEG: DWT provides multi-resolution time-frequency localisation, preserving sharp transient structure at fine scales while capturing low-frequency rhythms at coarse scales, a property well matched to sleep spindles and delta waves; STFT provides uniform-resolution spectral snapshots well suited to quantifying sustained power changes in canonical bands. Using both transforms in tandem is therefore not redundant but complementary: DWT encodes transient morphology that STFT blurs, and STFT provides stable band-power estimates that DWT decomposes unevenly across scales. From the DWT coefficients at each level, four statistics are computed: energy, log-energy, Shannon entropy, and variance. From the STFT, relative power is computed for Delta (0.5--4~Hz), Theta (4--8~Hz), Alpha (8--13~Hz), Beta (13--30~Hz), and Gamma (30--45~Hz) bands, together with spectral entropy and spectral edge frequency at 95 percent. Ratio features including delta-to-beta, theta-to-alpha, and slow-to-fast power ratios encode clinically established spectral relationships without inflating the feature dimension. Rather than retaining the full time-frequency spectrogram, these handcrafted statistics capture the physiologically meaningful dimensions of the frequency domain in a 34-dimensional vector that is both compact and interpretable. This compactness reduces model complexity in the spectral branch, prevents overfitting on the smaller Sleep-EDF corpus, and, critically, allows each spectral feature to be traced back to a specific physiological interpretation during explainability analysis.

% For the downstream disorder detection task, which operates on subject-level rather than epoch-level inputs, a separate clinical feature extraction pipeline was applied to MESA PSG recordings. This pipeline exploits not only the raw sensor signals but also the output of the sleep staging module, using the predicted hypnogram to derive architecture-level statistics that would otherwise require manual scoring. Sleep architecture features, computed from the staged hypnogram, include total sleep time, time in bed, sleep onset latency, wake after sleep onset, and sleep efficiency, together with stage-specific durations, proportions, and transition metrics such as total stage transitions, transition rate per hour, and stage transition entropy. These architecture features are important because insomnia and restless legs syndrome manifest primarily through disruptions to sleep continuity and stage distribution rather than through focal respiratory or movement events. Respiratory features follow AASM event detection criteria applied to the airflow channel: apneas are detected as amplitude reductions of 90 percent or greater lasting at least 10 seconds, hypopneas as reductions of 30 percent or greater with associated desaturation, and events are classified as obstructive or central based on respiratory effort amplitude. Derived metrics include the apnea-hypopnea index, obstructive and central apnea indices, REM-specific AHI, and mean and maximum apnea duration, all of which are critical for apnea severity stratification. Oxygenation features from SpO2 analysis include mean saturation, nadir, T90 and T88 thresholds, desaturation indices ODI3 and ODI4, and desaturation area-under-curve, capturing the integrated oxygen burden associated with disordered breathing. Cardiac and HRV features computed from Pan-Tompkins R-peak detection encode autonomic nervous system regulation during sleep, which is disrupted in obstructive apnea through repetitive arousal and sympathetic activation. Periodic limb movement features derived from leg EMG following AASM criteria provide direct physiological markers for restless legs syndrome: PLM count, PLMI, burst rate, and inter-movement interval statistics encode the stereotyped leg movement pattern that distinguishes RLS from non-specific movement during sleep. REM behaviour disorder features from chin EMG quantify tonic and phasic muscle activity during REM sleep, providing a physiological signature for REM sleep disorders that co-occur with multiple movement and neurodegenerative conditions. The final clinical feature matrix contains 136 features per subject, after preprocessing steps including median imputation for HRV zero-values arising from failed R-peak detection, exclusion of near-zero variance features, log transformation of EEG power features to normalise heavy-tailed distributions, and pruning of highly correlated feature pairs ($|r| > 0.95$) to reduce multicollinearity.

% \begin{figure*}
%     \centering
%     \includegraphics[width=1\linewidth]{assets/prj48.png}
%     \caption{Core Technological methodology diagram os the end-to-end system.}
%     \label{fig:methodology}
% \end{figure*}

% \subsection{Sleep Staging Architecture}

% The sleep staging module implements a temporal-spectral fusion architecture that processes both the raw EEG waveform and the pre-computed spectral feature vector for each epoch, models sequential dependencies across epochs using a transformer encoder, and classifies each epoch into one of five sleep stages. The design builds upon converging evidence in recent literature that single-modality architectures face complementary failure modes~\cite{cong2024bits, fu2023temporal}, and that transformer-based sequence models outperform recurrent alternatives for capturing the long-range temporal dependencies inherent in sleep stage transitions~\cite{phan2022sleeptransformer}.

% The temporal encoder employs an Adaptive Atrous Pyramid~\cite{sun2025adg} adapted for one-dimensional physiological signals. The pyramid comprises four parallel convolutional branches operating with dilation rates of 1, 2, 4, and 8, enabling simultaneous receptive fields spanning approximately 10~ms, 20~ms, 40~ms, and 80~ms at 100~Hz sampling. This multi-scale processing is motivated by the multi-temporal nature of sleep EEG: brief K-complexes (200--800~ms) and sleep spindles (0.5--2~seconds) demand fine temporal resolution, while slow wave oscillations and the sustained spectral characteristics of N3 sleep require broader receptive fields. Critically, multi-scale dilation achieves this without pooling operations that would irreversibly discard temporal resolution, an important consideration for preserving the morphological precision required by the GradCAM explainability module. Following the MultiSEss architecture of Huang et al.~\cite{huang2025multisess}, each convolutional branch incorporates a Squeeze-and-Excitation (SE) block with reduction ratio 16, which performs global average pooling across the time dimension to compute per-channel statistics, then applies a two-layer bottleneck network with sigmoid activation to produce channel-wise attention weights. This channel attention mechanism learns to emphasise frequency-selective filters that are task-relevant for sleep staging while suppressing filters that respond to broadband artifact, providing a form of implicit signal quality weighting that is particularly valuable for heterogeneous clinical recordings. Branch outputs are aggregated through a learned adaptive gating mechanism with softmax-normalised weights, enabling dynamic emphasis of temporal scales that are most informative for specific sleep stages rather than applying fixed pooling across all scales. An adaptive average pooling layer then reduces the temporal dimension from 3,000 samples to a fixed 128-dimensional embedding, providing a compact and stage-discriminative representation of each epoch for downstream sequence modelling. The spectral encoder processes the 34-dimensional spectral feature vector through a shallow two-layer MLP with intermediate dimension 128, reflecting the principle that spectral features are already explicitly compressed representations of frequency content and do not require the deep hierarchical processing applied to raw waveforms.

% Three fusion strategies were evaluated to combine temporal and spectral embeddings. Concatenation serves as the simplest baseline, concatenating the two 128-dimensional vectors and projecting back to 128 dimensions through a linear layer. Gated fusion computes per-dimension soft blending weights via a sigmoid-activated gate network that receives both embeddings as input, producing a continuous interpolation between the modalities for each feature dimension independently. Bidirectional cross-attention fusion implements full transformer attention between the modalities, allowing temporal embeddings to query spectral key-value pairs and spectral embeddings to query temporal key-value pairs simultaneously, enabling the model to explicitly learn inter-modal correspondences such as the relationship between spindle morphology in the temporal stream and sigma-band power elevation in the spectral stream. Ablation experiments presented in Section~\ref{sec:results} demonstrated that simple concatenation achieves the highest staging accuracy among the three strategies, outperforming cross-attention fusion by 0.8 percentage points. This result is consistent with the interpretation that the complementary information in the two streams benefits from full preservation without gating-induced attenuation: concatenation retains both 128-dimensional representations without suppressing any modality-specific information, whereas gating and cross-attention introduce learned compression that may discard features relevant to specific sleep stages. Additionally, concatenation substantially reduces the computational cost and memory requirements of the fusion block compared to cross-attention, which requires quadratic operations over the combined embedding space, making the concatenation-based design more practical for deployment on resource-constrained clinical hardware.

% The sequence model component adopts a transformer encoder architecture to capture the temporal dependencies between epochs that are essential for accurate sleep staging. Prior work has explored a range of sequence models for this problem, including long short-term memory networks~\cite{phan2022sleeptransformer}, gated recurrent units, bidirectional LSTM and GRU variants, graph convolutional networks, temporal convolutional networks, and recurrent neural networks. Transformer encoders with self-attention have consistently emerged as the strongest approach in recent comparative evaluations, a finding attributed to their ability to capture non-local temporal context without the distance-dependent forgetting that limits recurrent models on long sequences. Sleep staging is inherently a temporally structured task: physiological sleep cycles of approximately 90 minutes produce systematic regularities in stage sequencing, with N3 predominating in the first half of the night and REM consolidating in the second, and with strong short-range dependencies governing permissible stage transitions. Modelling these dependencies requires contextual memory over sequences of hundreds of epochs, a regime where transformer global self-attention is particularly well suited because every epoch can attend directly to every other epoch in the window without information passing through intervening states, avoiding the gradient attenuation that impairs recurrent architectures on long sequences. The transformer encoder comprises six layers with four attention heads, embedding dimension 128, and feed-forward dimension 512, following the pre-normalization convention for improved training stability in deeper stacks. Sinusoidal positional encodings inject epoch index information, enabling the model to represent typical sleep architecture progression from early high-proportion N3 to late REM-dominant cycles. Global self-attention, in which every query attends to all keys in the window simultaneously, was selected over sparse or sliding-window attention variants: given the window size of 256 epochs and the 128-dimensional embedding, full attention is computationally tractable and avoids the locality bias of windowed attention that could impair modelling of the 90-minute ultradian cycle. The transformer output is passed to a linear classification head projecting from 128 to 5 dimensions, producing per-epoch logits over the five sleep stage classes. Auxiliary classification heads attached to the temporal and spectral encoder outputs independently provide per-modality predictions that support attribution of model decisions to specific input streams during explainability analysis and enable separate validation of each encoder's contribution during ablation studies.

% \subsection{Sleep Disorder Detection}

% The disorder detection module identifies three clinically prevalent sleep disorders from PSG recordings: obstructive sleep apnea, insomnia, and restless legs syndrome. This task is formulated as multi-label binary classification, with independent predictions made for each disorder per subject, recognising that the conditions frequently co-occur and share overlapping physiological signatures. Labels derive from clinically adjudicated questionnaire responses in the MESA dataset, where participants reported whether they had received each diagnosis. Positive class prevalence is 7.6 percent for apnea, 6.3 percent for insomnia, and 4.5 percent for restless legs syndrome, producing a severe class imbalance that precludes accuracy as a meaningful evaluation metric and motivates F1 score as the primary criterion, with precision-recall area under the curve as the secondary metric.

% Multiple classifier architectures were selected to span the interpretability-capacity spectrum while maintaining the explainability requirements of clinical deployment. Logistic regression with ElasticNet regularisation provides a linear baseline whose coefficients are directly interpretable as log-odds contributions of individual features, offering strong face validity for clinical review. Random forests provide non-linear feature interaction modelling through bootstrap-aggregated decision trees with 500 estimators, offering inherent feature importance estimates via mean decrease in impurity and robustness to outliers through majority voting across diverse trees, properties well suited to the heterogeneous clinical feature set. XGBoost employs gradient-boosted trees with second-order Taylor expansion of the loss function and regularised leaf weights, combining strong predictive performance with the ability to produce exact SHAP attributions through the TreeExplainer algorithm, which is central to the patient-facing feature importance explanations. Linear support vector machines with hinge loss provide a maximum-margin linear classifier that is computationally efficient and well calibrated to high-dimensional feature spaces after dimensionality reduction, though sensitivity to scale requires robust standardisation. Multilayer perceptrons implemented in PyTorch with architecture 256-128-64-1 capture non-linear feature interactions beyond the reach of tree ensembles and linear models, at the cost of requiring additional regularisation to avoid overfitting on the limited positive-class samples.

% All models were trained with configurations specifically adapted to the class imbalance. Binary cross-entropy loss employed positive class weighting equal to the ratio of negative to positive samples, which upweights gradient contributions from the minority class without altering the decision boundary. For the MLP, the AdamW optimiser with weight decay $10^{-4}$ and learning rate $10^{-4}$ was used, with early stopping on validation loss with patience of 10 epochs and restoration of the best-validation checkpoint to prevent overfitting. Classification thresholds were tuned independently per disorder using out-of-fold probabilities from training, selecting the threshold that maximises F1 score via precision-recall curve analysis, rather than defaulting to 0.5. This threshold tuning is essential under extreme class imbalance, where the optimal decision boundary typically lies substantially below 0.5 to recover sufficient recall for the minority positive class. Feature selection preceded classification within each cross-validation fold to prevent information leakage: a two-stage pipeline first applied mutual information filtering with threshold 0.002 to remove features with negligible marginal association with the target, reducing from 136 to approximately 80--100 features per disorder, then applied ElasticNet logistic regression with L1 ratio 0.5 and bootstrap stability selection over 30 resampled iterations, retaining only features selected in more than 30 percent of bootstrap runs. The ElasticNet penalty balances L1 sparsity, which enforces explicit feature selection, with L2 ridge regularisation, which encourages correlated features within physiological groups to remain together rather than arbitrarily discarding all but one representative. This group retention property is particularly important for insomnia, where no single dominant feature predicts the self-reported label but combinations of sleep continuity metrics collectively encode the disorder pattern. Model evaluation employed stratified 5-fold cross-validation maintaining class distribution within each fold, with performance metrics aggregated via macro-averaging across folds. Strict separation between feature selection and evaluation was maintained throughout, with the complete feature selection pipeline applied independently to each training fold.

% \subsection{Explainability Methods}

% The explainability pipeline produces interpretable outputs for two clinically distinct audiences: sleep specialists and clinicians who require technically precise, signal-level evidence to audit model predictions and make informed clinical decisions, and patients or general users who require accessible, jargon-free information to understand their condition, track progress, and engage with treatment. This two-tier design is motivated by the recognition that a single explanation format cannot simultaneously serve the technical validation needs of expert reviewers and the comprehension needs of lay recipients. The pipeline operates on outputs of both the sleep staging module and the disorder classification module, generating attributions, uncertainty estimates, and visualisations at both the epoch level and the subject level.

% \subsubsection{Clinical Tools for Specialist Review}

% Attention rollout produces per-epoch influence scores quantifying which epochs in the analysis window contributed most to the model's prediction for a given target epoch. Sleep staging is inherently context-dependent: the interpretation of an isolated N1 epoch differs depending on whether it is preceded by consolidated N2 or follows a REM period, and the transformer encoder explicitly models these dependencies through cross-epoch attention. Attention rollout makes these dependencies interpretable by accumulating attention weights across all transformer layers and normalising the resulting matrix to produce scalar influence scores per epoch. The implementation bypasses PyTorch's fast-path optimisation, which discards attention weights for memory efficiency, by monkey-patching each self-attention forward pass to force weight return. Resulting influence score profiles allow clinicians to identify epochs where the model's decision was driven by temporally distant or potentially spurious context, flagging cases where attention drift may indicate unreliable predictions. The inference is provided by the transformer encoder module and is computed jointly with the forward pass during model evaluation without requiring additional model calls.

% EpochGradCAM generates within-epoch signal localisation maps that highlight which segments of the 30-second EEG waveform drove the model's classification decision. Clinicians require signal-level evidence to verify that model decisions align with known EEG patterns, specifically whether the model attends to sleep spindles in N2, slow waves in N3, sawtooth waves in REM, or K-complexes at sleep onset, rather than to electrical artifacts or non-physiological signal perturbations. GradCAM achieves this by computing the gradient of the class score with respect to the feature maps of the AdaptiveAtrousPyramid projection layer, weighting the feature map activations by their gradient magnitudes to produce a spatially localised importance map, and upsampling the result to the 3,000-sample input resolution. The implementation registers forward hooks on the projection layer to capture feature maps during inference, then performs a single targeted backward pass with \texttt{retain\_grad()} to preserve intermediate gradients. Regions exceeding the 80th percentile of the smoothed importance map are highlighted as primary attention regions, with the top three regions annotated in the visualisation. This threshold-based region selection provides clinicians with a concise, interpretable summary of the three most influential signal segments rather than a continuous heatmap that requires subjective interpretation. The method is provided by the temporal encoder module and requires only a single additional backward pass per epoch being explained.

% Monte Carlo Dropout uncertainty quantification provides per-epoch predictive entropy estimates by approximating the Bayesian posterior over model parameters through stochastic inference. In clinical AI systems, distinguishing high-confidence correct predictions from uncertain predictions that coincide with correct labels by chance is essential for calibrated trust: a model may produce a correct staging decision for the wrong reasons when uncertainty is high, and a clinician who sees only the prediction without uncertainty information cannot distinguish reliable classifications from unreliable ones. The implementation forces all dropout layers to training mode during evaluation via explicit traversal of the module tree, then performs 30 stochastic forward passes with independent dropout masks, producing a distribution of class probability vectors. Predictive entropy is computed as $H = -\sum_{k} \bar{p}_k \log \bar{p}_k$ where $\bar{p}_k$ is the mean probability of class $k$ across the 30 samples. High entropy indicates disagreement among the sampled model realisations and thus low confidence in the prediction, while low entropy indicates consistent predictions across realisations regardless of their direction. Confidence is reported to clinicians using a three-tier label mapped from entropy ratios normalised by the maximum possible entropy for five classes ($\log 5$): ratios below 0.45 are reported as High confidence, ratios between 0.45 and 0.72 as Medium confidence, and ratios above 0.72 as Low confidence, providing a clinically interpretable uncertainty summary without requiring familiarity with information-theoretic quantities.

% The automated review flagging system synthesises the uncertainty and prediction disagreement signals into a prioritised review queue that directs clinician attention to the most uncertain and potentially erroneous predictions. This system addresses the practical constraint that clinicians cannot manually review every epoch of a full-night recording, which may contain 900 or more epochs per patient. Epochs are flagged when either of two conditions holds: the model prediction disagrees with the reference label, or the predictive entropy exceeds the 75th percentile of the recording's entropy distribution. The percentile-based threshold adapts to subject-specific baseline uncertainty, recognising that recordings from subjects with unusual sleep architecture or electrode placement issues exhibit systematically higher baseline entropy throughout the recording, and that a fixed entropy threshold would over-flag such recordings or under-flag easier recordings. Flagged epochs are presented in a timeline visualisation with sleep-stage colour coding, per-epoch confidence indicators, and annotation arrows highlighting the five highest-priority epochs ordered by combined entropy and prediction disagreement score. This prioritised presentation enables systematic clinical review within a practical time budget.

% \subsubsection{Patient-Facing Explanations}

% Patient-facing explanations translate technical model outputs into accessible information that supports informed self-management, treatment adherence, and productive clinical consultations. All features in this tier are designed to communicate through familiar concepts and avoid technical terminology including entropy, gradient, attention weight, or model confidence, replacing these with clinically grounded analogues expressed in plain language.

% The disorder risk score communicates model output as a percentage risk estimate for each of the three disorders, mapped to categorical risk bands using clinically informed thresholds. Converting predicted probabilities to risk percentages preserves the ordinal information of the probabilistic output while presenting it in a format that patients can relate to personal health risk communication conventions. Risk bands (low: below 30 percent, moderate: 30--60 percent, high: above 60 percent) provide discrete categories that support decision-making without overwhelming patients with numeric precision that the model's calibration does not support. These scores are produced by the disorder classification module and formatted for patient consumption in the presentation layer, enabling patients to track risk trajectories across successive recordings and observe the impact of lifestyle changes on their disorder profile.

% Feature importance visualisation presents the top contributing physiological factors for each disorder prediction as percentage influence scores derived from SHAP values computed via the TreeExplainer algorithm on the XGBoost classifier, or via signed permutation importance for other classifier types. Each contributing feature is displayed with its patient-readable label, the patient's recorded value, and the population normal range, contextualising the patient's measurement within clinically meaningful reference bounds. For example, the system may indicate that wake after sleep onset contributed 32 percent of the insomnia risk score, that the patient's recorded value of 65 minutes substantially exceeds the normal reference of below 30 minutes, and that this metric reflects the total time spent awake after initially falling asleep. This multi-layered presentation enables patients to connect an abstract contribution percentage to a specific sleep behaviour and its physiological significance, supporting the kind of concrete self-monitoring goals that improve treatment adherence.

% Counterfactual explanations present hypothetical scenarios identifying what minimal changes to the patient's sleep profile would shift the disorder prediction to a lower risk category. Counterfactuals address a fundamental limitation of feature importance explanations, which describe what drove the current prediction but do not indicate what the patient could change to improve their outcome. The implementation identifies the minimum-perturbation feature change that crosses the classification threshold for each disorder, using the disorder-specific feature importances to prioritise physiologically actionable features over those outside the patient's control. Crucially, counterfactual descriptions are expressed in natural clinical language rather than as numerical feature deltas: instead of reporting that reducing WASO\textsubscript{sec} by 2160 would lower the predicted probability, the system states that reducing time spent awake during the night from 65 minutes to approximately 30 minutes would lower insomnia risk from the high to the moderate band. This translation from feature-space mathematics to natural-language clinical guidance is achieved through a rule-based description layer that maps each feature identifier and its direction of change to a pre-authored plain-language template grounded in evidence-based sleep hygiene recommendations.

% The AI advisor chat provides an interactive, conversational interface through which patients can ask follow-up questions about their results, explore the implications of specific findings, and receive personalised guidance on sleep hygiene and lifestyle modifications. A large language model is deployed with a system prompt that injects the patient's specific feature values, risk scores, and counterfactual scenarios as structured context, grounds model responses in the patient's actual results rather than generic sleep advice, and explicitly prohibits technical jargon while requiring plain language appropriate for lay audiences. The conversational format allows patients to seek clarification on unfamiliar terms, explore hypothetical questions not covered by the static explanation panels, and receive empathetic explanations that address the emotional dimension of receiving a risk assessment. The advisor operates as a complementary channel to the structured explanations rather than a replacement, with the structured panels providing persistent reference information and the chat interface providing on-demand interactive clarification.



% \section{METHODS}
% The proposed framework integrates a unified end-to-end pipeline for automated sleep analysis, disorder detection, and personalised explanation. The system is structured into five modular stages: dataset acquisition, signal preprocessing, sleep staging architecture design, multi-label disorder classification, and a dual-tier explainability pipeline. This modularity ensures that each component can be independently validated while contributing to a cohesive diagnostic workflow. The overall architecture, illustrating the flow from raw polysomnography (PSG) signals to actionable patient recommendations, is depicted in Figure~\ref{fig:methodology}.

% \subsection{Dataset Description}

% This study utilises two large-scale public polysomnography corpora to evaluate staging performance and clinical relevance. The Sleep-EDF Expanded dataset~\cite{sleepedfx}, sourced from PhysioNet~\cite{physiobank}, provides 197 whole-night recordings from healthy and sleep-disturbed adults, used here for supervised training and validation of the staging model. Complementing this, the Multi-Ethnic Study of Atherosclerosis (MESA) Sleep dataset~\cite{mesa}, accessed via the National Sleep Research Resource~\cite{nsrr}, provides over 2,200 recordings from a diverse adult cohort. While Sleep-EDF serves as the staging benchmark, MESA is used for disorder association analysis due to its clinically adjudicated labels for apnea, insomnia, and restless legs syndrome. Together, these datasets enable a robust assessment of model generalisation across different recording environments and population demographics.
% \subsection{Data Preprocessing and Representation}
% To ensure consistency across datasets and reduce computational overhead, only the Fpz-Cz EEG derivation was retained for sleep staging. Raw EDF recordings were processed using the MNE-Python library and segmented into non-overlapping 30-second epochs, resulting in 3,000 samples per epoch at a 100~Hz sampling rate. This sampling rate was selected to satisfy the Nyquist criterion for all clinically relevant bands up to 45~Hz while preserving the morphology of transient events like K-complexes and sleep spindles. Non-physiological labels, such as Movement Time and Unknown, were excluded, and recordings were trimmed to remove extended wake periods at the boundaries. The processed stage transitions for sample subjects are illustrated in Fig.~\ref{fig:stagetransition}.
% The system employs a dual-stream representation to capture both raw signal morphology and frequency-domain characteristics. The temporal stream provides the raw EEG waveform, while the spectral stream consists of a 34-dimensional vector of pre-computed statistics. For the spectral stream, Discrete Wavelet Transform (DWT) using the Daubechies-4 basis and Short-Time Fourier Transform (STFT) were used to extract energy, entropy, and relative power across Delta, Theta, Alpha, Beta, and Gamma bands. This combination is motivated by the need to preserve both sharp transient structures via DWT and stable band-power estimates via STFT. This dual-representation prevents the information loss associated with pure spectral compression while avoiding the redundancy of learning frequency filters from scratch.
% For subject-level disorder detection, a separate clinical feature extraction pipeline was applied to the MESA dataset. This pipeline integrates raw sensor data with architecture-level statistics derived from the predicted hypnogram, including total sleep time, sleep efficiency, and stage-specific durations. Respiratory features were extracted following AASM criteria for apnea and hypopnea, while cardiac and HRV features were computed using Pan-Tompkins R-peak detection. Periodic limb movement markers and REM-behaviour features were derived from leg and chin EMG signals, respectively. The final 136-feature matrix underwent median imputation for HRV zeros, log-transformation of EEG power, and pruning of highly correlated pairs ($|r| > 0.95$) to ensure numerical stability.
% \subsection{Sleep Staging Architecture}
% The sleep staging module implements a temporal-spectral fusion transformer that processes raw EEG waveforms and spectral features simultaneously to classify epochs into five stages. The temporal encoder utilises an Adaptive Atrous Pyramid~\cite{sun2025adg} with four parallel convolutional branches using dilation rates of 1, 2, 4, and 8 to capture multi-scale temporal patterns. Each branch incorporates a Squeeze-and-Excitation (SE) block to dynamically emphasise task-relevant filters while suppressing broadband artifacts. Branch outputs are aggregated via a learned adaptive gating mechanism and pooled into a 128-dimensional embedding. Simultaneously, the spectral encoder processes the 34-dimensional feature vector through a two-layer MLP to produce a matching 128-dimensional representation.
% These two embeddings are combined using a concatenation fusion strategy, which was found to outperform gated and cross-attention mechanisms by preserving all modality-specific information. The fused representation is then passed to a transformer encoder comprising six layers with four attention heads and sinusoidal positional encodings. The transformer models long-range temporal dependencies across a window of 256 epochs, effectively capturing the 90-minute ultradian sleep cycles. The final output is processed by a linear classification head to produce per-epoch logits for the five sleep stages. Auxiliary heads are attached to the individual encoders to enable separate validation of temporal and spectral contributions.
% \subsection{Sleep Disorder Detection}
% Disorder detection is formulated as a multi-label binary classification task for obstructive sleep apnea, insomnia, and restless legs syndrome. Due to the severe class imbalance (4--9\% positive rate), F1 score is used as the primary metric, with thresholds tuned via out-of-fold probabilities to maximise recall for minority classes. A suite of classifiers was evaluated, including ElasticNet Logistic Regression for linear interpretability, Random Forests for non-linear interaction modelling, and PyTorch MLPs (256-128-64-1) for high-capacity learning. All models utilised a positive class weighting in the binary cross-entropy loss to prevent the classifiers from collapsing to a majority-negative prediction.
% To manage the high dimensionality of the clinical feature set, a two-stage feature selection pipeline was implemented. First, a mutual information filter with a threshold of 0.002 removed features with negligible marginal association, reducing the feature space to approximately 80--100 per disorder. Second, ElasticNet bootstrap stability selection was applied over 30 resampled iterations, retaining only features selected in more than 30\% of the runs. This L1+L2 approach was specifically chosen to preserve complementary feature groups, which proved critical for improving insomnia detection performance. Model evaluation was conducted using stratified 5-fold cross-validation to maintain class distributions across all folds.
% \subsection{Explainability Methods}
% The explainability pipeline provides a two-tier system of interpretations tailored for clinicians and patients. For specialists, attention rollout accumulates weights across all transformer layers to quantify the temporal influence of adjacent epochs on a given prediction. Signal-level evidence is provided via EpochGradCAM, which computes the gradient of the class score relative to the AdaptiveAtrousPyramid feature maps to localise influential signal segments within the 30-second window. Furthermore, Monte Carlo (MC) Dropout approximates the Bayesian posterior by performing 30 stochastic forward passes to compute predictive entropy. Epochs are flagged for manual review if the entropy exceeds the 75th percentile of the recording or if the prediction disagrees with the ground truth.
% Patient-facing explanations translate these technical outputs into accessible, jargon-free guidance. Predicted probabilities are converted into percentage risk scores and mapped to low, moderate, or high risk bands. Feature importance is visualised using SHAP values from the XGBoost classifier, presenting the top contributing factors alongside the patient's actual values and population normal ranges. Counterfactual explanations identify the minimum feature perturbations required to shift a prediction to a lower risk category, expressed in natural clinical language. Finally, an AI advisor chat leverages a large language model to answer follow-up questions, grounding its responses in the patient's specific results while strictly avoiding technical terminology.


\section{METHODS}
The proposed framework integrates a unified end-to-end pipeline for automated sleep analysis, disorder detection, and personalised explanation, structured into five modular stages. The first stage encompasses dataset acquisition and description, followed by a signal preprocessing pipeline that transforms raw data into model-ready representations. The third stage details the sleep staging architecture, focusing on the temporal-spectral fusion transformer. The fourth stage describes the multi-label sleep disorder detection system, and the final stage presents the dual-tier explainability pipeline. This modular organization allows for independent validation of each component while ensuring a seamless flow from raw polysomnography (PSG) signals to clinically actionable patient recommendations.

\subsection{Dataset Description}
The landscape of automated sleep research is supported by numerous publicly available polysomnography corpora, including the seminal Sleep-EDF, the cardiovascular-focused SHHS, the older male population-centric MrOS, and the family-based CFS studies. Other specialized resources such as the WSC, APPLES, UCDDB, MIT PhysioNet, MIBIH, and CAP datasets provide diverse benchmarks for sleep staging and disorder detection, though they vary significantly in acquisition protocols and demographics. This work utilizes the Sleep-EDF Expanded dataset for the supervised training and validation of the sleep staging model and the Multi-Ethnic Study of Atherosclerosis (MESA) Sleep dataset for disorder association analysis. This combination enables the assessment of staging performance on a well-characterized benchmark while ensuring clinical relevance through the adjudicated labels of a large, population-representative cohort.

\subsubsection{Sleep EDF Expanded}
The Sleep-EDF Expanded dataset~\cite{sleepedfx}, sourced via PhysioNet~\cite{physiobank}, consists of 197 whole-night PSG recordings from healthy and sleep-disturbed adult subjects. The cohort spans a wide age range, approximately 25 to 101 years, and includes recordings from both the Sleep-Cassette (SC) and Sleep-Telemetry (ST) subsets. Data were acquired using laboratory-grade equipment, with the Fpz--Cz EEG channel sampled at 100~Hz to provide the primary signal for sleep staging. Annotations follow the Rechtschaffen and Kales standard, segmenting the recordings into 30-second epochs across five stages: Wake, N1, N2, N3, and REM. The resulting distribution reflects the natural class imbalance of nocturnal sleep architecture, necessitating the use of specialized loss functions to ensure robust detection of minority stages.

\subsubsection{MESA}
The MESA Sleep dataset~\cite{mesa}, accessed through the National Sleep Research Resource (NSRR)~\cite{nsrr}, comprises over 2,200 recordings from a diverse adult cohort aged approximately 45--84 years. Unlike the staging-focused Sleep-EDF, MESA is a population-based study designed to investigate the relationship between sleep characteristics and cardiometabolic health. PSG acquisition followed standardized protocols, including multi-channel EEG and respiratory signals, with disorder-level labels for obstructive sleep apnea, insomnia, and restless legs syndrome. These labels are derived from clinically adjudicated self-reported questionnaires, providing a target for disorder association analysis. This dataset is primarily used for generalisation testing and disorder-level evaluation, bridging the gap between automated staging and clinical diagnosis.

\subsection{Data Preprocessing and Representation}
The preprocessing pipeline transforms raw polysomnography (PSG) recordings into three distinct representations: raw temporal tensors, spectral feature vectors, and subject-level clinical matrices. The first stage converts raw EDF files into tensors using the MNE-Python library. To ensure cross-dataset consistency, only the Fpz-Cz EEG channel is retained, and all signals are resampled to 100~Hz to satisfy the Nyquist criterion for all clinically relevant bands up to 45~Hz. Signals are segmented into non-overlapping 30-second epochs, resulting in 3,000 samples per epoch, with non-physiological labels such as Movement Time and Unknown excluded to ensure data purity.

From these temporal tensors, a 34-dimensional spectral feature vector is extracted to capture sustained frequency-domain characteristics. This is achieved through a dual-transform approach: a Discrete Wavelet Transform (DWT) using the Daubechies-4 basis at decomposition level 5, and a Short-Time Fourier Transform (STFT) with a 2-second Hann window. DWT is used to extract energy, log-energy, Shannon entropy, and variance to capture transient morphological events, while STFT provides relative power estimates for Delta, Theta, Alpha, Beta, and Gamma bands. Together with spectral entropy and edge frequency, these features provide a noise-robust summary of the epoch's spectral content.

Finally, a subject-level clinical feature matrix is constructed by aggregating raw sensor signals and the output of the upstream sleep staging module. Architecture-level statistics, such as total sleep time, sleep efficiency, and stage durations, are derived from the predicted hypnogram. Respiratory and oxygenation features are extracted following AASM criteria, and cardiac HRV features are computed via Pan-Tompkins R-peak detection. Leg EMG signals are processed to identify periodic limb movement markers. The resulting 136-feature matrix is refined via median imputation for HRV zeros, log-transformation of EEG power, and pruning of highly correlated pairs ($\|r\| > 0.95$) to reduce multicollinearity.

\subsection{Sleep Staging Model Design}
Accurate sleep staging is a prerequisite for identifying sleep disorders, as architectural abnormalities serve as primary diagnostic indicators. The proposed sleep staging architecture consists of four primary modules: a multi-scale feature extractor, a modality fusion module, a sequence model, and a linear classification head. The system takes a dual-stream input consisting of a raw EEG waveform (3,000 samples) and a 34-dimensional spectral vector, and outputs per-epoch logits across five sleep stages. This design ensures that the model captures both the transient morphological patterns and the sustained spectral rhythms essential for precise staging.

The feature extractor utilizes an Adaptive Atrous Pyramid~\cite{sun2025adg}, which employs four parallel convolutional branches with dilation rates of 1, 2, 4, and 8. This design is inspired by the need to capture multi-temporal patterns, ranging from brief K-complexes to prolonged slow-wave oscillations, without using pooling operations that would discard temporal resolution. Each branch incorporates a Squeeze-and-Excitation (SE) block, following the design principles of BiTS-SleepNet, to dynamically emphasize task-relevant filters while suppressing broadband artifacts. Branch outputs are aggregated through a learned adaptive gating mechanism and pooled into a 128-dimensional embedding.

The fusion module combines the temporal and spectral embeddings using a concatenation strategy. While gated fusion and bidirectional cross-attention (inspired by TSANet) were experimentally evaluated, concatenation was found to be the most effective approach by preserving all modality-specific information without attenuation. This allows the model to explicitly leverage the complementary relationship between raw signal morphology and band-power estimates. The fused 256-dimensional representation is projected back to 128 dimensions to maintain a compact embedding space for the subsequent sequence model.

The sequence module employs a Transformer encoder to model long-range temporal dependencies across a window of 256 epochs. This architecture is inspired by the Sleep Transformer and Explainable Vision Transformer, utilizing self-attention to capture non-local context without the distance-dependent forgetting associated with recurrent networks. Sinusoidal positional encodings are used to represent the typical progression of sleep cycles, such as the transition from N3-dominant early sleep to REM-dominant late sleep. Six layers with four attention heads are used to ensure a high-capacity representation of the nocturnal sleep architecture.

The final classification head consists of a linear layer projecting the 128-dimensional transformer output to 5 dimensions. This head produces the final logits for the five sleep stages (Wake, N1, N2, N3, REM). To support explainability, auxiliary heads are attached to the individual temporal and spectral encoders, enabling the system to attribute predictions to specific input streams. This design ensures that the model remains computationally efficient while providing a robust mapping from multi-modal features to clinical stages.

\subsection{Sleep Disorder Detection}
Automated disorder detection is formulated as a multi-label binary classification task for obstructive sleep apnea, insomnia, and restless legs syndrome. While deep learning offers high capacity, traditional machine learning models often provide better interpretability and stability for subject-level diagnosis given the limited number of positive samples. Therefore, this work employs a suite of classifiers, including ElasticNet Logistic Regression, Random Forests, and XGBoost, alongside a PyTorch MLP (256-128-64-1) to span the interpretability-capacity spectrum. This hybrid approach ensures that the system can capture both linear clinical correlates and non-linear physiological interactions.

The process replicates clinical diagnosis by formulating clinical variables based on AASM guidelines. Features are derived from the processed PSG signals and the predicted hypnogram, focusing on respiratory indices (AHI, ODI), sleep continuity metrics (WASO, Sleep Efficiency), and movement markers (PLMI). By structuring the feature set to mirror clinical indicators, the model's decision process remains aligned with medical standards. This formulation ensures that the resulting predictions are grounded in established physiological markers rather than spurious correlations.

Technical AI settings are optimized to handle the severe class imbalance (4--9\% positive rate) present in the MESA dataset. All models use a positive class weighting in the binary cross-entropy loss to prevent collapse to majority-negative predictions. A two-stage feature selection pipeline is implemented, using a mutual information filter (threshold 0.002) followed by ElasticNet bootstrap stability selection over 30 iterations. This L1+L2 approach preserves complementary feature groups, which was critical for improving insomnia detection. Model performance is validated using stratified 5-fold cross-validation with thresholds tuned via out-of-fold probabilities to maximize F1 scores.

\subsection{Explainability Features}
The explainability pipeline provides a two-tier system of interpretations tailored to the distinct needs of clinicians and patients. The need for explainability in sleep medicine is driven by the requirement for clinical validation, the need to identify model uncertainty, and the goal of improving patient engagement in treatment. By translating black-box model outputs into human-interpretable evidence, the system bridges the gap between high-accuracy deep learning and practical clinical utility. The pipeline serves a dual user base: sleep technologists who require signal-level evidence for auditing, and patients who require accessible guidance for self-management.

\subsubsection{Clinician's Features}
For clinicians, the system provides three technical tools for model auditing. Attention rollout accumulates self-attention weights across all transformer layers to produce a per-epoch influence score, quantifying which epochs in the window drove each staging prediction. EpochGradCAM computes the gradient of the predicted class score with respect to the output of the \texttt{AdaptiveAtrousPyramid} projection layer (\texttt{self.proj}), producing a 3,000-sample importance map that is smoothed with a 25-sample moving average and normalised between the 5th and 95th percentiles. The resulting heatmap is rendered using the jet colourmap, with warm regions indicating the signal segments that most influenced the classification decision. Monte Carlo (MC) Dropout approximates the Bayesian posterior by performing 30 stochastic forward passes with all dropout layers forced to training mode. The output confidence ratio is derived from the predictive entropy normalised by $\log(5)$, the maximum possible entropy for five classes. Confidence tiers are calibrated on the full Sleep-EDF dataset: epochs are assigned High (green), Medium (orange), or Low (red) confidence based on data-driven thresholds that target a 10\% high-error rate. Epochs are flagged for manual review when the prediction disagrees with the ground truth or when the confidence ratio exceeds the 75th percentile of the recording.

Figure~\ref{fig:exp_dashboard} shows the integrated explainability dashboard for a sample recording, combining the hypnogram comparison, attention rollout influence scores, MC-Dropout uncertainty, and GradCAM signal localisation for a flagged epoch. Figure~\ref{fig:gradcam_main} shows GradCAM heatmaps overlaid on raw EEG for three sample epochs, illustrating how the model localises attention to physiologically meaningful signal regions. Figure~\ref{fig:mc_confidence} shows the per-epoch MC-Dropout confidence distribution across three patients, with epochs coloured by confidence tier.

\begin{figure*}[t]
    \centering
    \includegraphics[width=\linewidth]{exp.png}
    \caption{Integrated explainability output for a sample recording. Top panel: ground truth hypnogram (solid) versus model predictions (dashed) across 256 epochs. Second panel: per-epoch attention rollout influence scores. Third panel: MC-Dropout predictive entropy per epoch coloured by uncertainty level. Bottom panel: GradCAM localisation for epoch 117 (predicted N2, true REM, uncertainty 1.113), with red regions indicating the signal segments that most influenced the misclassification.}
    \label{fig:exp_dashboard}
\end{figure*}

\begin{figure}[t]
    \centering
    \includegraphics[width=\linewidth]{signal importance_clinician.png}
    \caption{GradCAM signal importance for three sample epochs (C2 output). Each panel overlays the raw EEG signal (black) with a GradCAM heatmap rendered using the jet colourmap, where warm (red/yellow) regions indicate the signal segments that most influenced the model's staging decision. The heatmap is computed from gradients of the predicted class score with respect to the \texttt{AdaptiveAtrousPyramid} projection layer, smoothed with a 25-sample moving average and normalised between the 5th and 95th percentiles.}
    \label{fig:gradcam_main}
\end{figure}

\begin{figure}[t]
    \centering
    \includegraphics[width=\linewidth]{confidence calliberated_clinician.png}
    \caption{MC-Dropout confidence strip across three patients (C3 output). Each row shows a per-epoch confidence timeline coloured by tier: High (green), Medium (orange), Low (red). Confidence tiers are derived from predictive entropy normalised by $\log(5)$, with thresholds calibrated on the full Sleep-EDF dataset to target a 10\% high-error rate. Patient 2 shows two low-confidence epochs near stage transitions, which are automatically flagged for manual review.}
    \label{fig:mc_confidence}
\end{figure}

\subsubsection{Patient's features}
Patient-facing explanations translate technical outputs into accessible, jargon-free guidance. Predicted probabilities from the XGBoost classifier are converted into percentage risk scores and mapped to three risk bands: Low (below 30\%), Moderate (30--60\%), and High (above 60\%). Feature importance is visualised using SHAP values computed via the \texttt{TreeExplainer} algorithm on the XGBoost model, presenting the top contributing factors as signed bar charts alongside population normal ranges to contextualise the patient's health status. Counterfactual explanations identify the minimum required percentage reduction in each feature to shift the prediction toward a lower risk band, expressed as actionable lifestyle targets with difficulty ratings (Easy, Medium, Hard). A conversational AI advisor uses the Groq Cloud API (Llama model) with a system prompt that injects the patient's SHAP-derived risk context in plain English, strictly prohibiting technical jargon including entropy, SHAP, model accuracy, and coefficient values.

Figure~\ref{fig:shap_main} shows the SHAP feature importance comparison for a control subject versus a high-risk apnea subject, illustrating how the model's decision is dominated by CPAP device use for the positive case. Figure~\ref{fig:counterfactual_main} shows the counterfactual analysis for the same high-risk apnea subject, presenting the minimum required reductions in each feature to shift the prediction to a lower risk band.

\begin{figure}[t]
    \centering
    \includegraphics[width=\linewidth]{Feature Importance_clinicians.png}
    \caption{SHAP feature importance for apnea classification comparing a control subject (risk = 0\%, left) and a high-risk subject (risk = 99\%, right). For the control subject, CPAP device use and AHI are the strongest risk-reducing factors. For the high-risk subject, CPAP device use dominates with a SHAP value of 5.73, confirming that the self-reported label is strongly driven by prior diagnosis and treatment history rather than raw physiological severity.}
    \label{fig:shap_main}
\end{figure}

\begin{figure}[t]
    \centering
    \includegraphics[width=\linewidth]{Counterfactuals feature_patient.png}
    \caption{Counterfactual analysis for a high-risk apnea subject (risk = 99\%). Each bar shows the minimum required percentage reduction in the corresponding feature to shift the prediction toward a lower risk band. CPAP device use requires a 70\% reduction, while all other features require 30\% reductions. Difficulty is colour-coded: Easy (green), Medium (orange), Hard (red).}
    \label{fig:counterfactual_main}
\end{figure}


\section{RESULTS}\label{sec:results}

\subsection{System Setup}

All experiments were conducted on Kaggle GPU kernels using dual NVIDIA Tesla T4 GPUs (30 GB VRAM, 15 GB RAM). The temporal-spectral fusion model for sleep staging was implemented in PyTorch 2.0, while disorder classification utilized scikit-learn 1.3, XGBoost 2.0, and PyTorch. Python 3.10 served as the primary environment, with pandas and NumPy for data manipulation. GPU utilization remained between 85--95\% during training.

The MESA Sleep Study dataset was obtained via the NSRR, and the Sleep-EDF Expanded dataset via PhysioNet. These provided 2,056 subjects with complete PSG and questionnaire data for disorder detection, and 197 subjects from Sleep-EDF for staging evaluation. Raw EEG signals were sampled at 100 Hz from the F4-M1 channel. Spectral features were pre-computed across delta, theta, alpha, sigma, and beta bands. The dataset was split into training and test sets using an 80/20 stratified split (random seed 42) at the subject level to prevent data leakage.

\subsection{Experimental Design}

The sleep staging experiments included a comprehensive ablation study comparing temporal-only, spectral-only, and fusion architectures. Table~\ref{tab:ablation} presents the results of this ablation analysis. The temporal-only model using raw EEG signals achieved a macro F1 of 0.1906, with strong N2 detection (F1 0.5335) but poor performance on N1, N3, and REM stages due to the limited feature representation. The spectral-only model utilizing hand-crafted frequency domain features achieved significantly improved macro F1 of 0.58, with balanced performance across all stages including wake (F1 0.78), N2 (F1 0.74), and N3 (F1 0.68). This established that spectral features carry essential discriminative information for sleep staging. Three fusion mechanisms were evaluated: simple concatenation, gated fusion, and bidirectional cross-attention. All fusion variants substantially outperformed single-modality approaches, with the concatenation fusion achieving the highest overall accuracy of 79.91 percent and macro F1 of 0.7877.

\begin{table}
\centering
\caption{Ablation study comparing temporal-only, spectral-only, and fusion architectures}
\label{tab:ablation}
\begin{tabular}{lcccccc}
\toprule
Architecture & Wake & N1 & N2 & N3 & REM & MF1 \\
\midrule
Temporal Only & 0.420 & 0.000 & 0.534 & 0.000 & 0.000 & 0.191 \\
Spectral Only & 0.780 & 0.250 & 0.740 & 0.680 & 0.460 & 0.582 \\
Cross-Attention & 0.911 & 0.588 & 0.754 & 0.758 & 0.889 & 0.780 \\
Gated Fusion & 0.911 & 0.600 & 0.763 & 0.748 & 0.893 & 0.783 \\
Concatenation & 0.917 & 0.600 & 0.765 & 0.768 & 0.889 & 0.788 \\
\bottomrule
\end{tabular}
\end{table}

Three primary experiments were designed to evaluate disorder detection performance under different feature selection and data source configurations. Experiment 3 established a baseline using L1 regularization for feature selection on questionnaire data. Experiment 4 replaced L1 with ElasticNet stability selection to test whether the L2 component would preserve complementary feature groups. Experiment 5 evaluated PSG signal-derived features independently to assess the alignment between physiological signals and self-reported disorder labels.

Evaluation metrics were selected to address the severe class imbalance present in all three disorder targets. Positive case prevalence ranged from 4.5 percent for restless legs syndrome to 8.9 percent for sleep apnea. F1 score served as the primary evaluation metric, with precision-recall area under the curve providing secondary validation. ROC-AUC was reported for completeness but not emphasized, as it can be misleadingly optimistic at low positive rates. Classification thresholds were tuned using out-of-fold probabilities on the training set to maximize F1 score, rather than using the default 0.5 threshold which would produce near-zero recall for minority classes.

The feature selection pipeline for questionnaire data in Experiments 3 and 4 followed a multi-stage approach. Initial preprocessing removed features with greater than 40 percent missing values, followed by variance filtering to remove constant features. Correlation pruning eliminated one feature from each pair with absolute Pearson correlation exceeding 0.95. Mutual information ranking and XGBoost feature importance provided intermediate guidance, with stability selection serving as the final feature selection mechanism. Experiment 3 used L1-regularized logistic regression with 30 bootstrap iterations, while Experiment 4 replaced this with ElasticNet using an L1 ratio of 0.5 to balance sparsity with group retention. Experiment 5 applied a hard mutual information threshold of 0.003 for PSG signal features, as percentile-based thresholds resolved to zero on this dataset.

\subsection{Performance Evaluation}

The temporal-spectral fusion model achieved strong sleep staging performance across five classes. Table~\ref{tab:sota_comparison} presents the comparison with state-of-the-art methods on Sleep-EDF Expanded dataset. The proposed concatenation fusion approach achieved a weighted average accuracy of 81.04 percent with particularly strong performance on N1 detection at 59.95 percent F1, surpassing all compared methods including BiTS-SleepNet (50.23 percent), TSA-Net (30.11 percent), and Sleep Transformer (48.5 percent). The cross-attention fusion variant achieved 80.26 percent accuracy with N1 F1 of 58.82 percent, while the gated concatenation variant achieved 80.67 percent accuracy with N1 F1 of 59.98 percent. Wake detection remained consistently strong across all fusion variants, ranging from 91.10 percent to 91.67 percent F1. The superior N1 performance is clinically significant as this transitional stage has historically posed the greatest challenge for automated sleep staging systems.

\begin{table*}
\centering
\caption{Comparison with state-of-the-art sleep staging methods on Sleep-EDF Expanded}
\label{tab:sota_comparison}
\adjustbox{max width=\textwidth}{
\begin{tabular}{lccccccc}
\toprule
\multirow{2}{*}{Method} &
\multicolumn{5}{c}{Per-Class F1 (\%)} &
\multirow{2}{*}{Accuracy (\%)} \\
\cmidrule(lr){2-6}
 & N1 & N2 & N3 & REM & Wake & \\
\midrule
MultiSEss & 47.62 & 83.51 & 75.36 & 75.11 & 91.58 & 82.30 \\
ADG SleepNet & 53.70 & 90.00 & 85.50 & 84.80 & 91.30 & 85.10 \\
BiTS-SleepNet & 50.23 & 87.05 & 82.50 & 85.33 & 93.33 & 85.09 \\
TSA-Net & 30.11 & 84.94 & 80.02 & 81.03 & 91.71 & 82.21 \\
Sleep Transformer & 48.50 & 86.50 & 80.90 & 84.60 & 93.50 & 84.90 \\
\midrule
Cross Attention Fusion (Ours) & 58.82 & 75.43 & 75.84 & 88.86 & 91.10 & 80.26 \\
Gated Concatenation (Ours) & 59.98 & 76.27 & 74.79 & 89.25 & 91.13 & 80.67 \\
Concatenation Fusion (Ours) & 59.95 & 76.54 & 76.78 & 88.90 & 91.67 & 81.04 \\
\bottomrule
\end{tabular}
}
\end{table*}

Figure~\ref{fig:CM_stage} presents the confusion matrix for the temporal-spectral concatenation fusion model. The normalized confusion matrix reveals strong diagonal classification for wake, N2, N3, and REM stages, with wake showing 91.67 percent accuracy and N3 demonstrating 76.78 percent accuracy. The confusion patterns align with clinical expectations, with N1 epochs showing distributed misclassification across adjacent stages due to the transitional spectral characteristics of this stage.

\begin{figure*}
    \centering
    \includegraphics[width=1\linewidth]{assets/confusion_matrix.png}
    \caption{Normalized confusion matrix for the temporal-spectral concatenation fusion model.}
    \label{fig:CM_stage}
\end{figure*}

Table~\ref{tab:disorder_results} summarizes the disorder detection performance across Experiments 3, 4, and 5. Experiment 4 achieved the best overall questionnaire-based results, with sleep apnea F1 of 0.7843, insomnia F1 of 0.3881, and RLS F1 of 0.1553. The ElasticNet approach improved insomnia detection by 35.9 percent relative to Experiment 3, confirming that the L2 component successfully preserved complementary feature groups that pure L1 regularization discarded. XGBoost demonstrated the greatest robustness to high-dimensional feature spaces, achieving the highest apnea F1 score of 0.7843 with AUROC of 0.8835 and PR-AUC of 0.7646. The MLP implemented in PyTorch achieved the best insomnia F1 of 0.3881 with AUROC of 0.7597, suggesting that non-linear interactions between sleep architecture features benefit from deep learning architectures. However, all models struggled with RLS detection from questionnaire features, with F1 scores remaining below 0.16.

\begin{table*}
\centering
\caption{Disorder detection performance across experiments}
\label{tab:disorder_results}
\adjustbox{max width=\textwidth}{
\begin{tabular}{@{}lllccccccccc@{}}
\toprule
\multirow{2}{*}{Exp} & \multirow{2}{*}{Features} & \multirow{2}{*}{Selection} &
\multicolumn{3}{c}{Apnea} & \multicolumn{3}{c}{Insomnia} & \multicolumn{3}{c}{RLS} \\
\cmidrule(lr){4-6} \cmidrule(lr){7-9} \cmidrule(lr){10-12}
 & & & F1 & AUROC & PR-AUC & F1 & AUROC & PR-AUC & F1 & AUROC & PR-AUC \\
\midrule
3 & Questionnaire & L1 & 0.782 & -- & -- & 0.286 & -- & -- & 0.166 & -- & -- \\
4 & Questionnaire & ElasticNet & \textbf{0.784} & 0.884 & 0.765 & \textbf{0.388} & 0.760 & 0.194 & 0.155 & 0.631 & 0.093 \\
5 & PSG Signals & MI threshold & 0.183 & 0.623 & 0.103 & 0.130 & 0.494 & 0.076 & \textbf{0.182} & 0.660 & 0.097 \\
\bottomrule
\end{tabular}
}
\end{table*}

Experiment 5 revealed a critical finding regarding label-feature alignment. PSG signal features achieved only an F1 of 0.183 for apnea detection despite using direct physiological measurements including AHI and oxygen desaturation indices. This counterintuitive result occurs because the disorder labels represent self-reported diagnoses rather than clinically measured severity. The questionnaire labels encode whether a person was told they have a disorder, which correlates with healthcare seeking behavior, CPAP usage history, and symptom awareness rather than raw physiological values. The one exception was RLS, where PSG features achieved F1 of 0.182, outperforming questionnaire-based detection. This reflects the presence of periodic limb movement markers in the PSG signals, which serve as direct physiological signatures of RLS.

The ablation study comparing L1 and ElasticNet regularization demonstrated the importance of regularization choice for multi-disorder detection. Pure L1 regularization in Experiment 3 selected only 9 stable features, all of which were strong standalone predictors of apnea including AHI, BMI, and CPAP usage history. This sparsity came at the cost of excluding features that were individually weak but jointly predictive of insomnia and RLS. ElasticNet in Experiment 4 retained 326 to 383 features per target, representing 52 to 61 percent of the original feature set. While this retention rate proved more permissive than ideal, it successfully captured complementary feature groups that L1 discarded. The 35.9 percent relative improvement in insomnia F1 demonstrates that sleep architecture features such as WASO, sleep efficiency, and stage percentages are individually noisy but collectively informative when combined through non-linear models. The mutual information pre-filter proved essential for merged datasets containing both questionnaire and PSG features, bringing the p/n ratio from 0.42 to below 0.15 and enabling meaningful feature selection even with 693 initial features.

\subsection{Discussion}

The experimental results reveal several important findings regarding multi-disorder sleep detection and the relationship between data sources, feature selection methods, and prediction targets. The ablation study demonstrates that neither temporal nor spectral features alone achieve optimal staging performance. The temporal-only model excelled at N2 detection but failed on N1, N3, and REM stages, suggesting that temporal patterns alone are insufficient for comprehensive staging. The spectral-only model achieved moderate performance across all stages, confirming that frequency domain features provide essential discriminative information. However, the fusion architectures combining both modalities substantially outperformed single-modality approaches, with the simple concatenation fusion achieving the highest accuracy. This suggests that temporal and spectral features capture complementary aspects of sleep EEG, and their combination enables robust classification across all five stages.

The most significant insight concerns the misalignment between self-reported disorder labels and raw physiological signals. Sleep apnea and insomnia labels in the MESA dataset represent patient-reported diagnoses rather than objectively measured severity. This distinction explains why questionnaire features containing behavioral correlates of healthcare utilization outperformed PSG signal features containing direct physiological measurements. A person with moderate AHI who has discussed symptoms with a physician receives a positive label, while a person with identical physiology who has not sought medical attention receives a negative label. This label definition fundamentally limits the ceiling performance of any model trained on physiological features alone.

The exception to this pattern was restless legs syndrome, where PSG-derived periodic limb movement features achieved the best performance. Unlike apnea and insomnia, RLS produces distinctive physiological signatures in leg EMG activity and associated arousal patterns that are captured by PLM counts and indices. This finding suggests that disorder-specific physiological markers can overcome label misalignment when such markers exist, but behavioral questionnaire features remain essential for conditions lacking clear physiological signatures.

The staging results demonstrate notable achievements in N1 classification, with the concatenation fusion achieving 59.95 percent F1, substantially exceeding specialized architectures such as BiTS-SleepNet (50.23 percent), TSA-Net (30.11 percent), and Sleep Transformer (48.5 percent). This performance is clinically significant because N1 represents a transitional state with spectral characteristics overlapping adjacent stages, historically posing the greatest challenge for automated sleep staging systems. The superior N1 detection indicates that the temporal-spectral fusion architecture effectively captures subtle spectral transitions that single-modality approaches miss. Strong wake detection (91.67 percent F1) and REM detection (88.90 percent F1) further demonstrate the model's clinical utility, as accurate identification of these endpoints is essential for reliable sleep architecture analysis.

The explainability pipeline addresses a critical gap in clinical acceptance of automated sleep analysis. Attention rollout provides epoch-level influence scores indicating which adjacent epochs drove each prediction, supporting clinical review of borderline cases. GradCAM localization identifies specific time regions within the 30-second epoch where the model focused, enabling clinicians to validate predictions against visible EEG events such as sleep spindles or K-complexes, as illustrated in Figure~\ref{fig:gradcam_main}. Monte Carlo dropout entropy quantifies prediction uncertainty, flagging epochs requiring manual review when entropy exceeds the 75th percentile of the recording or when predictions disagree with ground truth labels, as shown in Figures~\ref{fig:exp_dashboard} and~\ref{fig:mc_confidence}. The patient-facing SHAP importance and counterfactual outputs (Figures~\ref{fig:shap_main} and~\ref{fig:counterfactual_main}) confirm that the model's apnea predictions are dominated by CPAP device use, directly reflecting the label-feature alignment finding.

Several limitations constrain the current findings and suggest directions for future research. The RLS detection performance remains insufficient for clinical deployment, with the best F1 score of 0.182 indicating that current features do not adequately capture RLS pathology. Creating physiological labels based on PLM thresholds rather than self-report might improve this performance. Class imbalance across all three disorders limits the effective sample size for minority class learning, with positive cases representing less than 9 percent of the dataset for all targets. The reliance on self-reported labels introduces measurement error that likely reduces achievable performance ceilings regardless of model sophistication.

Future work should explore multi-task learning architectures that jointly model all three disorders, as insomnia and RLS share sleep architecture correlates that separate single-task models cannot exploit. Direct SHAP value computation from XGBoost via TreeExplainer would provide more accurate feature importance estimates than the permutation importance approximations currently used. Cross-dataset evaluation on held-out populations such as SHHS or CHAT would assess generalization beyond the MESA cohort. Finally, a structured recommendation engine mapping specific feature deviations to validated lifestyle interventions would address the personalized recommendation component of the research objectives beyond the current LLM-based advisor chat.

\section{CONCLUSION}

This work presented a unified explainable deep learning framework for sleep disorder detection that bridges the gap between accurate automated analysis and clinical utility. The temporal-spectral fusion architecture achieved strong sleep staging performance with macro F1 of 0.788, notably achieving 59.95 percent F1 on the challenging N1 stage that has historically posed the greatest difficulty for automated systems. The ablation study demonstrated that fusion approaches substantially outperform single-modality models, with simple concatenation proving competitive against more complex attention mechanisms.

For disorder detection, the ElasticNet stability selection approach improved insomnia detection by 35.9 percent over traditional L1 regularization by preserving complementary feature groups that pure sparsity constraints discarded. The analysis revealed an important finding regarding label-feature alignment: self-reported disorder labels correlate with behavioral questionnaire features rather than raw physiological signals, with the exception of restless legs syndrome where periodic limb movement markers provide direct physiological signatures.

The explainability pipeline addresses a critical barrier to clinical adoption by providing tailored explanations for distinct audiences. Clinicians receive attention-based temporal explanations, GradCAM signal localization, and Monte Carlo dropout uncertainty quantification, enabling informed review of automated predictions. Patients receive accessible risk scores, feature importance visualizations, counterfactual explanations identifying actionable changes, and conversational AI guidance supporting self-management. This dual-interface design recognizes that clinical utility requires not only accuracy but also transparency and actionability for all stakeholders.

Several limitations suggest directions for future work. The modest RLS detection performance indicates need for physiological label definitions based on PLM thresholds rather than self-report. Class imbalance across all disorders constrains minority class learning. The reliance on questionnaire labels introduces measurement error that bounds achievable performance regardless of model sophistication. Future research should explore multi-task learning architectures jointly modeling all three disorders, direct SHAP computation for accurate feature importance, cross-dataset evaluation on held-out populations, and structured recommendation engines mapping specific deviations to validated lifestyle interventions.

The broader implications of this work extend beyond technical performance metrics. By demonstrating that explainable AI can achieve competitive accuracy while providing clinically actionable insights, the framework supports a vision of human-AI collaboration in sleep medicine where automated systems handle routine analysis while escalating uncertain cases to human experts. The modular architecture enables deployment across settings from sleep laboratories to home monitoring devices, potentially expanding access to sleep disorder diagnosis for underserved populations. As sleep disorder prevalence continues rising globally, such transparent and scalable approaches will be essential for addressing the growing gap between diagnostic need and clinical capacity.

\section*{ETHICAL APPROVAL}

The MESA Sleep study was approved by institutional review boards at all participating sites. All participants provided written informed consent. This secondary analysis was conducted using de-identified data from the National Sleep Research Resource (NSRR) repository. The Sleep-EDF Expanded dataset is publicly available under open research licenses.

\bibliographystyle{IEEEtran}
\bibliography{references}

\begin{IEEEbiography}[{\includegraphics[width=1in,height=1.25in,clip,keepaspectratio]{assets/author1.png}}]
{Girish Sundararajan} is currently working toward the B.Tech. degree in Computer Science with Artificial Intelligence from the School of Computing, Amrita Vishwa Vidyapeetham, Amritapuri, India. His research interests include AI for neurodiversity, disability studies, and sustainable development.
\end{IEEEbiography}

\begin{IEEEbiography}[{\includegraphics[width=1in,height=1.25in,clip,keepaspectratio]{assets/author2.jpeg}}]
{Harishankar Binu Nair} is currently working toward the B.Tech. degree in Artificial Intelligence from the School of Computing, Amrita Vishwa Vidyapeetham, Amritapuri, India. His research interests include natural language processing, robotics, and AI for disability studies.
\end{IEEEbiography}

\begin{IEEEbiography}[{\includegraphics[width=1in,height=1.25in,clip,keepaspectratio]{assets/author3.png}}]
{Simi Surendran }  received the B.Tech. degree in Computer Science and Engineering from Mahatma Gandhi University and the M.Tech. degree in Wireless Networks from Amrita Vishwa Vidyapeetham, India. She is currently a faculty member and research fellow with the Department of Computer Science, Amrita Vishwa Vidyapeetham. She is pursuing the Ph.D. degree in wireless networks. Her research interests include wireless networks, stochastic algorithms, and machine learning.
\end{IEEEbiography}

\end{document}
